\begin{center}
    {\fontsize{14pt}{16pt}\selectfont\textbf{ABSTRAK}\par}
    \vspace{0.8cm}
\end{center}

Perkembangan teknologi informasi memberikan dampak signifikan terhadap berbagai sektor industri, termasuk perbankan. Penerapan sistem pembelajaran berbasis digital di lingkungan perusahaan terbukti meningkatkan efektivitas pelatihan karyawan. PT. BPR Jatim (Perseroda) sebagai Badan Usaha Milik Daerah (BUMD) di Jawa Timur membutuhkan sistem pelatihan internal yang efektif dan fleksibel untuk meningkatkan kompetensi karyawannya. Sebelumnya, perusahaan menggunakan platform \textit{e-Learning} eksternal pihak ketiga yang bersifat generik, memiliki keterbatasan fitur, dan tidak dapat disesuaikan dengan kurikulum internal serta alur bisnis pelatihan PT. BPR Jatim (Perseroda). Oleh karena itu, dikembangkan sistem \textit{e-Learning} berbasis \textit{web} dan \textit{mobile} yang bertujuan menyediakan platform pelatihan yang dapat diakses karyawan kapan saja dan di mana saja, serta mempermudah divisi Pengembangan Sumber Daya Manusia (PSDM) dalam mengelola materi pelatihan, menyelenggarakan kuis evaluasi, dan memantau progres pembelajaran karyawan secara terpusat. Metode pengembangan yang digunakan adalah rancang bangun perangkat lunak yang meliputi analisis kebutuhan sistem, perancangan antarmuka pengguna menggunakan Figma, serta implementasi menggunakan Next.js sebagai \textit{framework web admin}, Flutter sebagai \textit{framework mobile} untuk karyawan, dan Supabase sebagai layanan \textit{backend} yang menyediakan basis data PostgreSQL, autentikasi, serta penyimpanan file. Pengujian fungsional dilakukan menggunakan metode \textit{black-box testing} untuk memverifikasi kesesuaian fitur dengan kebutuhan pengguna. Hasil kerja praktik menunjukkan bahwa sistem \textit{e-Learning} multiplatform yang dikembangkan berhasil memfasilitasi proses pembelajaran karyawan secara lebih terstruktur dan interaktif, mampu menyajikan laporan evaluasi secara \textit{real-time} kepada manajemen, serta memberikan fleksibilitas waktu pembelajaran tanpa mengganggu jam kerja operasional.

\vspace{1cm}
\noindent \textbf{Kata Kunci:} \textit{E-Learning}, Pelatihan Karyawan, \textit{Multiplatform}.

\newpage
